\documentclass[10pt]{report}

\title{Full steam ahead(ers), Captain!}
\author{Joshua Clark}

\usepackage[a4paper, left=1.25in, right=1.25in]{geometry}
\usepackage{amsmath}
\usepackage{amssymb}
\usepackage[citestyle=authoryear]{biblatex}
\usepackage{backnaur}
\usepackage{url}
\bibliography{references}

\usepackage{graphicx}
\usepackage[unicode=true,bookmarks=true,bookmarksnumbered=false,bookmarksopen=false,breaklinks=true,pdfborder={0 0 0},colorlinks=false]{hyperref}
\hypersetup{pdftitle={E-Mail Header Information}, pdfauthor={Joshua Clark}}

\allowdisplaybreaks
\begin{document}

\thispagestyle{empty} 

\begin{center}
\begin{minipage}[c]{0.75\linewidth}%
 \centering %University logo
\includegraphics[width=0.4\linewidth]{oxford} 

\vspace{2cm}
 %Thesis title 
{\uppercase{\Large{}E-Mail Header Information}{\Large \par}} \vspace{2cm}
%Author's name 
{\Large{}Joshua Clark}{\Large \par}

\vspace{2cm}
%Degree 
{\Large{}Fourth Year Project Report for the Final Honour School of Computer Science}{\Large \par}

 \vspace{2cm}
 %Date 
{\Large{}May 2016} %
\end{minipage}
\par\end{center}

\clearpage{}


\chapter*{Abstract}
After extensive public education, fewer people are now clicking on links in e-mails that are disguised as phishing attacks, though the threat still remains, and considerable amounts of work has gone into exploring the demographics most likely to be targeted.  As the number of technically literate people grows, this sort of attack is increasingly unlikely to be successful.  Therefore, malicious entities are more likely to attempt to attack people based on the information leaked in their emails, and more specifically, the header, which most people are less likely to have some degree of control over.  

This report discusses the existing research into the information leaked by e-mail headers and presents a tool to extract such information.

\chapter*{Acknowledgements}
I want to thank my supervisor, Jason Nurse for his assistance throughout the year; my tutor, Peter Jeavons, for his unfailing help and support throughout my time at Oxford.  I would like to acknowledge the support from my family and partner, Agata, for encouraging me.

\tableofcontents
\chapter{Introduction}
E-mail systems are now so integrated into our modern lives that we struggle to cope without them.  E-mails ubiquity is also one of its largest weaknesses, a fact recognised very early on.  The first spam email was sent in 1978, as documented by \cite{templeton}.  After spam came phishing, first described by \cite{felix1987system}, with the first-real world use being against the customers of America Online, an ISP.  However, this still relies on the targets providing their data for malicious purposes.  One of the first e-mail viruses to spread was the Happy99 virus, which, other than propagating itself, had no other effect on infected systems.  Later viruses would target credit-card and banking information.  However, all of these techniques rely on the malicious email being received and its contents being opened.  There are fewer instances recorded, however, of the information flow being sent the other way.  A more subtle attack will focus on the information being sent from a legitimate user to an attacker.  It is easy enough for an individual to read an e-mail header and identify interesting elements, however, on a large scale, this quickly becomes more difficult.

\section{Typical Use}

On starting the application, the user will provide an e-mail that they wish to have analysed.  This will then be parsed, and relevant information presented in a table and other pictoral based representations.

Lastly, an option is available to view the information about security vulnerabilities in a separate webpage.

\section{Document Conventions}

By convention, when class diagrams are used, ovals will represent
traits/abstract classes, with rectangles representing concrete implementations.
Objects are indicated using bold lines in the diagram, and are similar
in behaviour to statically declared objects in many languages.

Whenever \texttt{this font face} is used, the text is referring to
either some implementation level object or class, or text entered into the application or used for testing.

\chapter{Existing Research}
In \cite{nurse2015investigating}, the idea of using the information available in an email header was mooted, turning the previously standard threat of malware and phishing contained in received e-mails on its head, and instead presenting the threat in outgoing emails, and the personally identifying information (PII) contained therein.  Many emails leaked information about employers, e-mail services and applications used, and IP address.  Initial examination of a variety of e-mail headers found within my own inbox also revealed a plethora of information, including phone carriers, preferred languages, and system usernames.  It is conceivable therefore, that it is possible to automate at least part of this, and present the information that can be extracted, in a white-hat tool to allow people to audit the information that they are revealing.  The obvious malicious use-case involves using such information as part of a spear-phishing exercise.

An alternative vulnerability presents itself in the information about systems that may be revealed.  Many email clients embed identifying information, and there are multiple databases available to allow specific threats to be identified.  This could allow a malicious entity to compromise the security of a target machine, and gain access to the data stored on that machine and available on any connected network devices.  Work started in \cite{joshi2013extracting} discusses the need to aggregate data about vulnerabilities from multiple sources to present a more complete and coherent picture, which is also likely to then contain more accurate data.

\cite{Al-zarouni_tracinge-mail} presents an alternative set of results, describing how an individual can seek to protect themselves against malicious e-mails, using the contents of e-mail headers.  Various discrepancies between forged e-mail addressses and legitimate messages are described.
\chapter{Design}
\section{Requirements}
The program would be expected to satisfy the following minimal requirements
in order for it to be considered successful: 
\begin{description}
\item [{Accuracy}]  -- any information produced by the parser should be reflective of the input e-mail
\item [{Representation}] -- the produced visualisation should be intuitive to read: each element should be presented separately from the others, and clearly labelled. 
\item [{Portability}] -- the visual output produced by the program should be available to the user in a variety of formats. 
\item [{Interactivity}] the program should produce sensible warnings when
an e-mail that is not possible to parse has been entered. 
\end{description}


\chapter{Definitions}\label{chap:def}
The following covers the essential definitions required for the notation and concepts that will be discussed in this document.
\section{Parsing}
In order to aid the parsing of the e-mail header, a combination of regular expressions and context-free grammars are needed, and defined as follows.
\subsection{Alphabets and Languages}
A set of symbols, usually denoted as $\Sigma$.  A language is a subset of $\mathcal P (\Sigma)$.

The following special classes are provided as part of the Perl-Compatible Regular Expression library, and are subsets of the alphabet of Unicode characters, defined in \cite{php_group_gutmans_lerdorf_suraski_boerger}.
\begin{description}
\item[alnum] -- letters and digits
\item[alpha] -- letters
\item[ascii] -- the set of ASCII characters (character codes 0 --- 127)
\item[blank] -- tabs or blank spaces
\item[cntrl] -- control characters
\item[digit] -- decimal digits
\item[graph] -- printing characters (excluding spaces)
\item[lower] -- lower-case letters
\item[print] -- printing characters (including spaces)
\item[punct] -- punctuation marks (printing characters excluding letters and spaces)
\item[space] -- white space
\item[upper] -- upper case letters
\item[word] -- ``word'' characters (same 
\item[xdigit] -- hexadecimal digits
\end{description}
\subsection{Regular Languages}
Regular languages are defined as follows:
\begin{itemize}
\item $\emptyset$ and $\{\epsilon\}$ are regular languages
\item for each $a\in\Sigma$, $\{a\}$ is a regular language
\item if $A$ and $B$ are both regular, $A\cup B$, $A\cdot B$ and $A^*$ are regular languages.
\subitem $A\cup B$ is the union of two languages.  $A\cup B = \{s : s\in A \lor s \in B\}$
\subitem $A\cdot B$ is the concatentation of two languages.  $A\cdot B = \{ ab : a \in A, b \in B\}$
\subitem $A^*$ is the Kleene star of a language. 
\begin{align*}
    A_0&=\{\epsilon\}\\
    A_1&= A\\
    A_{i+1} &= \{ aa' : a \in A_i, a'\in A\}\\
    A^* &= \bigcup_{i\in\mathbb N} A_i
\end{align*}
\end{itemize}

\subsection{Context-Free Grammars}
A context-free grammar $G$ is defined as $G=\left(V,\Sigma, R,S\right)$ where:
\begin{itemize}
\item $V$ is a variable.
\item $\Sigma$ is the alphabet of symbols.
\item $R$ is a relation defined over $V\rightarrow \left(V\cup\Sigma\right)^*$
\item $S$ is the start symbol
\end{itemize}

For example, $\langle \text S \rangle$ is the field name with the associated productions $\langle \text T \rangle \, \langle \text U \rangle$, where $T$ and $U$ are productions.
\begin{bnf*}
	\bnfprod{S}{\bnfpn{T} \bnfsp{} \bnfpn{U}}
\end{bnf*}
For example, $\langle \text S \rangle$ is the field name with the associated productions $a \, \langle \text U \rangle$, where $a$ is a terminal symbol.
\begin{bnf*}
	\bnfprod{S}{\bnftd{a} \bnfsp{} \bnfpn{U}}
\end{bnf*}
This is then extended in the following ways used in the RFC syntax.

The square brackets are used to indicate an optional element.
\begin{bnf*}
\bnfprod{field}{\bnfpn{field-name} \bnfts{:} \bnfsp{}[ \bnfpn{field-body} ]\bnfsp{} \bnfts{CRLF}}\\
\end{bnf*}

The asterisk is used to indicate an element that appears 0 or more times. $n*$ is used to indicate a component that repeats $n$ or more times.
\begin{bnf*}
\bnfprod{fields}{\bnfpn{dates}\bnfsp \bnfpn{source} \bnfsp 1\!*\bnfpn{destination} \bnfsp * \bnfpn{optional-fields}}\\
\end{bnf*}

The hash-symbol is used to indicate an element that appears a certain number of times. $m*n$ is used to indicate a component that repeats at least $m$ times and at most $n$ times.
\begin{bnf*}
\bnfprod{fields}{\bnfpn{dates}\bnfsp \bnfpn{source} \bnfsp 1\!\#\bnfpn{destination} \bnfsp * \bnfpn{optional-fields}}\\
\end{bnf*}

The $|$ is used to indicate a selection between a pair of elements.
\begin{bnf*}
\bnfprod{fields}{\bnfts{a}\bnfor \bnfts{b}}\\
\end{bnf*}

\section{Database Queries}
The following notations will be used for the CVE database queries.
\subsection{Set-Theoretic Operators}
The operators $F\cup G$, $F \cap G$, $F\setminus G$ behave as is expected for these operators, resulting in the union, intersection and difference of the sets.  The only proviso being that the atrribute names must match.
\subsection{Selection}
\[\sigma_{\text{product}=\text{thunderbird}}D\]
The above notation is used to indicate a search over the attribute named ``product'' for the string ``thunderbird'' in the database table $D$.  As a single database is only being used, this may be occasionally elided.  The output of this function is another object of the same type as $D$.

\subsection{Projection}
\[\pi_{\text{product}}D\]
The above notation is used to indicate a projection on the attribute named ``product'' in the database table $D$.  The output of this function is another object of the same type as $D$.

\subsection{Composition}
The above functions results can be coposed repeatedly to produce more specific search queries.

\chapter{Implementation}\label{chap:imp}
\section{Overview}

The  analysis is implemented as a series of stages, firstly, the e-mail header
is parsed, to extract important information to a predefined set of Java objects.
This is followed by the analysis phase, where the resultant data is passed to a
set of analyser modules, each running separately.  Finally, this information is
presented to the user.  This chapter presents each of these stages in detail.

\section{Definitions}

The following covers the essential definitions required for the notation and
concepts that will be discussed in this document.

\subsection{Parsing}

In order to aid the parsing of the e-mail header, a combination of regular
expressions and context-free grammars are needed, and defined as follows.

\paragraph{Alphabets and Languages}

A set of symbols, usually denoted as $\Sigma$.  A language is a subset of
$\mathcal P (\Sigma)$.

The following special classes are provided as part of the Perl-Compatible
Regular Expression library, and are subsets of the alphabet of Unicode
characters, defined in~\cite{php_group_gutmans_lerdorf_suraski_boerger}.

\begin{description}

\item[alnum] --- letters and digits

\item[alpha] --- letters

\item[ascii] --- the set of ASCII characters (character codes 0 --- 127)

\item[blank] --- tabs or blank spaces

\item[cntrl] --- control characters

\item[digit] --- decimal digits

\item[graph] --- printing characters (excluding spaces)

\item[lower] --- lower-case letters

\item[print] --- printing characters (including spaces)

\item[punct] --- punctuation marks (printing characters excluding letters and spaces)

\item[space] --- white space

\item[upper] --- upper case letters

\item[word] --- ``word'' characters (same

\item[xdigit] --- hexadecimal digits

\end{description}
\subsubsection{Regular Languages}
Regular languages are defined as follows:
\begin{itemize}
\item $\emptyset$ and $\{\epsilon\}$ are regular languages
\item for each $a\in\Sigma$, $\{a\}$ is a regular language
\item if $A$ and $B$ are both regular, $A\cup B$, $A\cdot B$ and $A^*$ are regular languages.
\subitem{ $A\cup B$ is the union of two languages.  $A\cup B = \{s : s\in A \lor s \in B\}$}
\subitem{ $A\cdot B$ is the concatentation of two languages.  $A\cdot B = \{ ab : a \in A, b \in B\}$}
\subitem{ $A^*$ is the Kleene star of a language.}
\begin{align*}
    A_0&=\{\epsilon\}\\
    A_1&= A\\
    A_{i+1} &= \{ aa' : a \in A_i, a'\in A\}\\
    A^* &= \bigcup_{i\in\mathbb N} A_i
\end{align*}
\end{itemize}

\subsubsection{Context-Free Grammars}
A context-free grammar $G$ is defined as $G=\left(V,\Sigma, R,S\right)$ where:
\begin{itemize}
\item $V$ is a variable.
\item $\Sigma$ is the alphabet of symbols.
\item $R$ is a relation defined over $V\rightarrow \left(V\cup\Sigma\right)^*$
\item $S$ is the start symbol
\end{itemize}

For example, $\langle \text S \rangle$ is the field name with the associated
productions $\langle \text T \rangle \, \langle \text U \rangle$, where $T$ and
$U$ are productions.

\begin{bnf*}
	\bnfprod{S}{\bnfpn{T} \bnfsp{} \bnfpn{U}}
\end{bnf*}

For example, $\langle \text S \rangle$ is the field name with the associated
productions $a \, \langle \text U \rangle$, where $a$ is a terminal symbol.

\begin{bnf*}
	\bnfprod{S}{\bnftd{a} \bnfsp{} \bnfpn{U}}
\end{bnf*}
This is then extended in the following ways used in the RFC syntax.

The square brackets are used to indicate an optional element.
\begin{bnf*}
\bnfprod{field}{\bnfpn{field-name} \bnfts{:} \bnfsp{}[ \bnfpn{field-body} ]\bnfsp{} \bnfts{CRLF}}\\
\end{bnf*}

The asterisk is used to indicate an element that appears 0 or more times. $n*$
is used to indicate a component that repeats $n$ or more times.

\begin{bnf*}
\bnfprod{fields}{\bnfpn{dates}\bnfsp \bnfpn{source} \bnfsp 1\!*\bnfpn{destination} \bnfsp * \bnfpn{optional-fields}}\\
\end{bnf*}

The hash-symbol is used to indicate an element that appears a certain number of
times. $m*n$ is used to indicate a component that repeats at least $m$ times and
at most $n$ times.

\begin{bnf*}
\bnfprod{fields}{\bnfpn{dates}\bnfsp \bnfpn{source} \bnfsp 1\!\#\bnfpn{destination} \bnfsp * \bnfpn{optional-fields}}\\
\end{bnf*}

The $|$ is used to indicate a selection between a pair of elements.
\begin{bnf*}
\bnfprod{fields}{\bnfts{a}\bnfor \bnfts{b}}\\
\end{bnf*}

\subsection{Database Queries}
The following notations will be used for the CVE database queries.

\paragraph{Set-Theoretic Operators}
The operators $F\cup G$, $F \cap G$, $F\setminus G$ behave as is expected for
these operators, resulting in the union, intersection and difference of the
sets.  The only proviso being that the atrribute names must match.

\paragraph{Selection} \[\sigma_{\text{product}=\text{thunderbird}}D\]
The above notation is used to indicate a search over the attribute named
``product'' for the string ``thunderbird'' in the database table $D$.  As a
single database is only being used, this may be occasionally elided.  The output
of this function is another object of the same type as $D$.

\paragraph{Projection}
\[\pi_{\text{product}}D\]
The above notation is used to indicate a projection on the attribute named
``product'' in the database table $D$.  The output of this function is another
object of the same type as $D$.

\paragraph{Composition}
The above functions results can be coposed repeatedly to produce more specific
search queries.

\section{Data Extraction and Parsing}

The parser's operation completes in a number of stages, following RFC822
(\cite{RFC0822}).  The header is divided up into two disjoint sections, the
routing information (\texttt{Received from...}) and the key-value map of other
pertinent information.

\subsection{Received fields}

The received fields are the most complicated part of the e-mail header to parse,
as they are described by a non-trivial grammar, presented below.

\begin{bnf*}
\bnfprod{message}{\bnfpn{fields}\bnfsp *(\bnfts{CRLF} \bnfsp *\bnftd{text})}\\
\bnfprod{fields}{\bnfpn{dates}\bnfsp \bnfpn{source} \bnfsp 1\!*\bnfpn{destination} \bnfsp * \bnfpn{optional-fields}}\\
\bnfprod{field}{\bnfpn{field-name} \bnfts{:} \bnfsp [ \bnfpn{field-body} ]\bnfsp \bnfts{CRLF}}\\
\bnfprod{field-name}{\bnftd{any word consisting of CHAR, excluding CTLs, SPACE, and ``'':''}} \\
\bnfprod{field-body}{\bnfpn{field-body-contents} \bnfsp [\bnfts{CRLF} \bnfsp \bnftd{LWSP-char}\bnfsp  \bnfpn{field-body}]}\\
\bnfprod{field-body-contents}{\bnftd{ASCII characters}}\\
\bnfprod{source}{[\bnfpn{trace}] \bnfsp \bnfpn{originator} [\bnfpn{resent}]}\\
\bnfprod{trace}{\bnfpn{return}\bnfsp 1\!* \bnfpn{received}}\\
\bnfprod{return}{\bnfts{Return-path:}\bnfsp{} \bnfpn{route-addr}}\\
\bnfprod{recieved}{\bnfts{Received:}}\\
\bnfprod{cont.}{[\bnfts{from}\bnfsp\bnfpn{domain}]}\\
\bnfprod{cont.}{[\bnfts{by}\bnfsp\bnfpn{domain}]}\\
\bnfprod{cont.}{[\bnfts{via}\bnfsp\bnfpn{atom}]}\\
\bnfprod{cont.}{*(\bnfts{with}\bnfsp\bnfpn{atom})}\\
\bnfprod{cont.}{[\bnfts{id}\bnfsp\bnfpn{msg-id}]}\\
\bnfprod{cont.}{[\bnfts{for}\bnfsp\bnfpn{addr-spec}]}\\
\bnfprod{cont.}{\bnfts{;}\bnfsp\bnfpn{date-time}}\\
\bnfprod{msg-id}{\bnfts{$<$}\bnfpn{addr-spec}\bnfts{$>$}}\\
\bnfprod{addr-spec}{\bnfpn{local-part}\bnfsp\bnfts{@}\bnfsp\bnfpn{domain}}\\
\bnfprod{local-part}{\bnfpn{word}\bnfsp *(\bnfts{.}\bnfsp\bnfpn{word})}\\
\bnfprod{word}{\bnfpn{atom}\bnfor\bnfpn{quoted-string}}\\
\bnfprod{domain}{\bnfpn{sub-domain} *(\bnfts{.}\bnfpn{sub-domain})}\\
\bnfprod{sub-domain}{\bnfpn{domain-ref}\bnfor\bnfpn{domain-literal}}\\
\bnfprod{domain-ref}{\bnfpn{atom}}\\
\bnfprod{date-time}{[ \bnftd{day,} ] \bnfsp \bnftd{date}\bnfsp \bnftd{time}}\\
\bnfprod{atom}{1\!*\bnftd{any character excluding specials, SPACE and CTLs}}\\
\end{bnf*}

An example field is as follows:
\begin{verbatim}
Received: from relay12.mail.ox.ac.uk (129.67.1.163)
    by HUB05.ad.oak.ox.ac.uk (163.1.154.231)
    with Microsoft SMTP Server id 14.3.169.1;
    Sat, 14 Nov 2015 10:55:35 +0000
\end{verbatim}
\subsection{Other fields}

These are read by a Python script and output to \texttt{STDOUT} to be read by
the Java parser in a consistent format.  These are then loaded into a hashmap to
allow quick lookup.

\section{Analysis}

After completing the parsing of the field, it is then ready to be analysed for
different features.  All of the analysers implement the \texttt{HeaderAnalyser}
interface, requiring information about the header to be analysed, and the
currently running application.  All of these then implement the
\texttt{Runnable} interface, allowing the class to be run asynchronously.

\subsection{Text-Based}

The fields from the header are analysed in different modules, with searches
being performed for specific strings.  Of particular interest to Oxford Nexus
users is the ``X-Oxford-Username'' string, containing the username of the
individual that sent the message.  As confirming the username is a fairly
standard security procedure for an IT support technician, having access to this
information could allow a phisher in a later stage of an attack to increase
their credibility.

In some cases, the likely keys that are being searched for are known in advance,
and can then be checked against the hash-map of entries.

An example of this approach is for the specific check for an Oxford username:

\begin{algorithm}
	\KwIn{Header}
	\KwOut{Any Oxford-based username that is found}
	\If{\texttt{X-Oxford-Username} $\in$ Header.KvMap}{ \Return{Header.KvMap(X-Oxford-Username)}\; }
	\caption{Lookup based on a known key}
\end{algorithm}

Alternatively, we may be interested in properties of the keys, necessitating a search over the keys.

\begin{algorithm}
	\KwIn{Header}
	\KwOut{Any information relating to Microsoft Exchange that is found}
	\ForEach{Key $k\in $ Header.KvMap}{
		\If{$k$ starts-with \texttt{X-MS-Exchange}}{
			\Return{Header.KvMap($k$)}\;
		}
	}
	\caption{Lookup based on a key property}
\end{algorithm}

\subsection{Database Queries}

Using the results gathered from the text-based queries and analysis of the
received fields, relevant software configurations are extracted and queried
against results in the CVE database.  These are then parsed and collated in
preparation for displaying the outputs.

As more information is found, more details of products used will also become
available.  These are added asynchronously.

\begin{algorithm}
	\KwIn{Header product name $p$}
	\KwOut{CVE Entries}
	cve-list $\gets\emptyset$\;
	\ForEach{$s\in\sigma_{\text{vector}\neq\text{LOCAL}}\sigma_{\text{product}=p} D$}{
		cve-builder$\gets$blank cve\;
		cve-builder.id$\gets\pi_{\text{CVE-ID}} s$\;
		$\ldots$ -- extract other features\;
		cve-list $\gets$ cve-list${}\cup{}$make(cve-builder)\;
	}
	\Return{cve-list}\;
	\caption{Extracting CVE entries}
\end{algorithm}

\section{Visualising the Results}

Using a pre-existing template, the results from the e-mail analysis will be
presented in a temporary webpage, which can then be saved independently.  Other
than the referenced JavaScript libraries, the document requires no additional
information or database access, allowing it to be quickly shared.

\chapter{Results}

\chapter{Conclusions}

\printbibliography
\end{document}
